% !TeX root = document.tex
\section*{Висновки}
У ході виконання дипломної роботи було повністю досягнуто поставлену мету та успішно вирішено всі сформульовані завдання. За результатами проведеного дослідження та програмної розробки отримано такі основні результати:

\begin{enumerate}
	\item \textbf{Проаналізовано} існуючі методи просторового масштабування цифрових зображень та досліджено математичну природу виникнення візуальних артефактів. Доведено, що використання класичних рівномірних сіток для глобальної поліноміальної інтерполяції є причиною виникнення феноменів Рунге та Гіббса на різких контурах.
	\item \textbf{Адаптовано} математичну модель цифрового зображення шляхом переходу до нерівномірних сіток дискретизації за вузлами Чебишева другого роду. Це забезпечило оптимальне згущення вузлів на межах інтервалу та гарантувало критично важливу алгоритмічну властивість математичної вкладеності сіток при зміні масштабу.
	\item \textbf{Поширено} другу барицентричну інтерполяційну формулу на двовимірний випадок. Успішно виконано її математичну декомпозицію, що дозволило кардинально знизити асимптотичну обчислювальну складність алгоритму. Для ефективної обробки різких градієнтів та локалізації похибок додатково розроблено локальну модифікацію методу в обмеженому просторовому околі.
	\item \textbf{Створено} програмне забезпечення мовою C++ із сучасним графічним інтерфейсом користувача на базі фреймворку Qt. Завдяки інтеграції з бібліотекою комп'ютерного зору OpenCV та застосуванню багатопотоковості, реалізовано апаратно-прискорене виконання алгоритму, здатне обробляти растрові матриці високої роздільної здатності в режимі реального часу.
	\item \textbf{Проведено} автоматизоване тестування розроблених рішень. За результатами об'єктивного порівняльного аналізу за допомогою метрик якості PSNR та SSIM.
\end{enumerate}

\textbf{Загальний підсумок:} У результаті проведеної роботи створено новий, обчислювально конкурентоспроможний алгоритм просторового масштабування зображень. 
Отримані результати формують надійний базис та відкривають два потужні напрямки для подальшого вдосконалення розробленого методу:
\begin{itemize}
	\item \textbf{Оптимізація швидкодії:} виведена матрична форма барицентричної інтерполяції ($G = B \cdot F \cdot A^\top$) не лише ефективно працює на центральному процесорі, але й ідеально адаптована для перенесення обчислень на графічні прискорювачі (GPU). Використання технологій апаратного прискорення, таких як CUDA або OpenCL, дозволить обробляти растрові масиви надвисокої роздільної здатності та відеопотоки у реальному часі.
	\item \textbf{Оптимізація якості наближення:} існують глибокі теоретичні резерви для покращення візуальної якості за допомогою математичного апарату де ла Валле-Пуссена. Застосування сум або фільтрів де ла Валле-Пуссена до побудованих чебишовських інтерполянтів дозволить ще ефективніше придушувати залишкові високочастотні осциляції на різких градієнтах. Це забезпечить досягнення майже найкращого наближення, усуваючи артефакти Гіббса без штучного розмиття структурної різкості зображення.
\end{itemize}

Алгоритм є повністю готовим до подальшого розвитку та впровадження у сучасні системи комп'ютерного зору та цифрової обробки сигналів.